\section{Dinamikus világháló}
1993-ban az NCSA (National Center for Supercomputing Applications) kifejlesztette a CGI-t (Common Gateway Interface), ezzel lehetővé téve, hogy a webszerverek külső programokat futtassanak (például Perl-szkripteket), és menet közben oldalakat generáljanak \cite{robinson2004cgi}. Ez volt az első nagy választóvonal a statikus és a dinamikus weboldalak között. A szerveroldali nyelvek (PHP, Python, ASP, Java) adatbázissal párosítva működtették a korai dinamikus oldalakat (fórumokat, online üzleteket) \cite{mdnServerSideIntro,mdnWebStandardsModel}. Ezeket a szerver által futtatott oldalakat minden egyes kérésnél újra kellett tölteni, ám a 2000-es évek végén az AJAX széles körű ismeretének és használatának köszönhetően az oldalak frissülni tudtak anélkül, hogy teljesen újra kellett volna őket tölteni, így gördülékenyebb felhasználói felületet biztosítottak \cite{garrett2005ajax,mdnAjax}. Ugyanakkor ezeknek az oldalaknak az alapja továbbra is a kliens és a szerver közötti kérés–válasz modellt követte \cite{fielding2022httpSemantics}.
